% !TEX program = xelatex
\documentclass[11pt]{article}

\usepackage{fontspec}
\setmainfont{CMU Serif}

\usepackage[margin=1in]{geometry}
\usepackage{amsmath,amssymb}
\usepackage{hyperref}
\usepackage{microtype}
\usepackage{enumitem}

\setlength{\parindent}{0pt}
\setlength{\parskip}{0.6\baselineskip}

\title{Executive Summary: GRA Мета-обнулёнка (многоуровневая архитектура)}
\author{}
\date{}

\begin{document}
\maketitle

\textbf{DOI:} \texttt{10.5281/zenodo.18470588}

\section{Коротко (1 абзац)}

\emph{GRA Мета-обнулёнка} --- это многоуровневый (иерархический) метод согласования многодоменных систем через минимизацию ``пены'' (квадратичного штрафа за рассогласование) относительно иерархии целей/проекторов. На практике метод можно интерпретировать как \textbf{регуляризатор} или \textbf{добавку к функции потерь}, которая подавляет междоменную интерференцию (конфликтные/внедиагональные компоненты) и тем самым уменьшает структурные противоречия, нестабильность оптимизации и ``галлюцинации'' в мультизадачных когнитивных моделях.

\section{Постановка проблемы (Google-scale)}

Во многих системах уровня Google (LLM, многомодальные модели, рекомендации, оптимизация инфраструктуры, биомедицинская аналитика) возникает общий класс ошибок:
\begin{itemize}[leftmargin=*, itemsep=2pt]
  \item \textbf{междоменное рассогласование} (конфликт ограничений/целей между задачами, датасетами, контекстами);
  \item \textbf{нестабильность обучения/оптимизации} (конкурирующие градиенты, дрейф представлений);
  \item \textbf{шум/спуриозные корреляции} как ``артефакты интерпретации'';
  \item \textbf{дефицит формального критерия} для ``согласованной истины'' в многоцелевой системе.
\end{itemize}

GRA формализует это как наличие ненулевой \emph{пены} между подсистемами/доменами на разных уровнях абстракции.

\section{Формальная идея (минимум необходимых формул)}

\subsection{Иерархия уровней}

Система задаётся уровнями $l=0,1,\dots,K$ и множествами подсистем/индексов $\mathcal{I}_l$.

\subsection{Проекторы целей}

Каждому уровню соответствует цель $G_l$ и проектор $\mathcal{P}_{G_l}$ на пространство состояний, удовлетворяющих этой цели.

\subsection{Пена уровня}

Ключевой штраф (пена) уровня $l$:
\[
\Phi^{(l)}(\Psi^{(l)},G_l)
= \sum_{\substack{\mathbf{a},\mathbf{b}\in\mathcal{I}_l\\ \mathbf{a}\neq\mathbf{b}}}
\left|\left\langle \Psi^{(\mathbf{a})}\middle|\mathcal{P}_{G_l}\middle|\Psi^{(\mathbf{b})}\right\rangle\right|^2.
\]

Интуитивно: чем сильнее ``внедиагональная'' связь между различными доменами при проекции на общую цель, тем больше пены и тем менее согласована система.

\subsection{Глобальный функционал (как loss/regularizer)}

Глобальный критерий (с затухающими весами по уровням):
\[
J_{\text{multiverse}}(\mathbf{\Psi})
= \sum_{l=0}^{K} \Lambda_l \sum_{\mathbf{a}\in\mathcal{I}_l} J^{(l)}\bigl(\Psi^{(\mathbf{a})}\bigr),
\qquad \Lambda_l=\lambda_0\alpha^l,\; 0<\alpha<1.
\]

В ML-интерпретации: добавка $\sum_l \Lambda_l\,\Phi^{(l)}$ может быть включена в общую функцию потерь как \textbf{структурный регуляризатор} для многодоменных представлений.

\section{Ключевые выводы/гарантии (в формулировке для ресёрча)}

При (i) коммутативности проекторов, (ii) согласованности иерархии (факторизация проекторов по подсистемам), (iii) достаточной размерности пространства, существует состояние $\mathbf{\Psi}^*$, в котором
\[
\Phi^{(l)}(\Psi^{(l)*},G_l)=0 \quad \forall\, l=0,\dots,K.
\]

Интерпретация: достижима согласованная конфигурация без междоменной ``интерференции'' на каждом уровне абстракции.

\section{Куда применять внутри Google (конкретные гипотезы)}

\subsection{DeepMind / Gemini (``anti-hallucination'' regularizer)}

Гипотеза: добавка $\Phi$ на уровне доменов/инструментов/контекстов уменьшит конфликтные представления и стабилизирует ``reasoning'' под ограничениями.

\subsection{Cloud / Infrastructure (distributed control)}

Гипотеза: моделирование узлов/подсистем как уровней и оптимизация пены уменьшит конфликты планирования и латентности в распределённых решениях.

\subsection{Health/Calico (согласование биомаркеров)}

Гипотеза: представление наборов биомаркеров как доменов и минимизация междоменной пены позволит искать устойчивые ``неподвижные точки'' состояния здоровья.

\section{План внедрения (минимально реалистичный)}

\begin{enumerate}[leftmargin=*, itemsep=3pt]
  \item \textbf{Whitepaper (10--12 стр.)}: сопоставить $\Phi$ с multi-objective optimization, constrained learning, information geometry.
  \item \textbf{Прототип на JAX}: реализация штрафа $\Phi$ и базового цикла оптимизации (батч-агрегация по доменам).
  \item \textbf{Эксперименты}: (а) multi-task instruction-tuning; (б) tool-use/agentic evals; (в) устойчивость к dataset shift.
  \item \textbf{Метрики}: hallucination rate, constraint satisfaction, gradient conflict (cosine), stability/variance of loss.
\end{enumerate}

\section{Что нужно от автора (2 пункта)}

\begin{itemize}[leftmargin=*, itemsep=2pt]
  \item Зафиксировать целевые сценарии (LLM, инфраструктура, биомед) и уровни $K$ / разбиение на домены.
  \item Уточнить семантику целей $G_l$ (как задаются проекторы) для выбранного пилота.
\end{itemize}

\end{document}